% Table: Diagonal vs full-covariance observation likelihood
% Source: offline analysis on HiddenCourtesyMerge-Sim-cleanobs/steps_with_belief.csv
\begin{table}[htbp]
\centering
\caption{Offline covariance-structure ablation for the three active
  observation features on held-out test-split in-window steps.
  The canonical policy evaluation uses the diagonal product-of-Gaussians
  likelihood; a full-covariance Gaussian fitted on training steps improves
  offline filter accuracy, indicating that residual feature correlations are
  practically relevant even after dropping \emph{rs}.}
\label{tab:covariance_ablation}
\small
\begin{tabular}{lccccc}
\toprule
\textbf{Likelihood} & \textbf{Step acc.} & \textbf{Final acc.} &
\textbf{NLL} & \textbf{Brier} & \textbf{ECE} \\
\midrule
Diagonal Gaussian & 82.8\% & 92.2\% & 0.428 & 0.131 & 0.068 \\
Full covariance Gaussian & 86.1\% & 97.7\% & 0.296 & 0.093 & 0.053 \\
\bottomrule
\end{tabular}
\begin{tablenotes}
\small
\item Values use 7{,}275 held-out in-window steps from 346 test episodes.
  Full-covariance policy rollouts were not rerun; Table~\ref{tab:policy_summary}
  therefore remains the canonical closed-loop result.
\end{tablenotes}
\end{table}
